\documentclass[letterpaper]{article} 
\usepackage[utf8]{inputenc}
\linespread{0.85}
\usepackage[T1]{fontenc}
\usepackage{amsmath}
\usepackage{amsfonts}
\usepackage{amssymb}
\usepackage{array}
\usepackage{fontspec}
\usepackage{booktabs}
\usepackage{hyperref}
\usepackage[version=4]{mhchem}
\usepackage{stmaryrd}
\usepackage[dvipsnames]{xcolor}
\colorlet{LightRubineRed}{RubineRed!70}
\colorlet{Mycolor1}{green!10!orange}
\definecolor{Mycolor2}{HTML}{00F9DE}
\usepackage{graphicx}
\usepackage{amsmath}
\usepackage{graphicx}
\usepackage{capt-of}
\usepackage{lipsum}
\usepackage{fancyvrb}
\usepackage{tabularx}
\usepackage{listings}
\usepackage[export]{adjustbox}
\graphicspath{ {./images/} }
\usepackage[utf8]{inputenc}
\usepackage[english]{babel}
\usepackage{float}
\usepackage{lipsum}
\usepackage{graphicx}
\usepackage{float}
\usepackage[margin=0.7in]{geometry}
\usepackage{amsmath}
\usepackage{graphicx}
\usepackage{capt-of}
\usepackage{tcolorbox}
\usepackage{lipsum}
\usepackage{graphicx}
\usepackage{float}
\usepackage{listings}
\definecolor{deepred}{RGB}{139,0,0}
\usepackage{hyperref} 
\usepackage{xcolor} % For custom colors
\lstset{
	language=Python,                % Choose the language (e.g., Python, C, R)
	basicstyle=\ttfamily\small, % Font size and type
	keywordstyle=\color{blue},  % Keywords color
	commentstyle=\color{gray},  % Comments color
	stringstyle=\color{red},    % String color
	numbers=left,               % Line numbers
	numberstyle=\tiny\color{gray}, % Line number style
	stepnumber=1,               % Numbering step
	breaklines=true,            % Auto line break
	backgroundcolor=\color{black!5}, % Light gray background
	frame=single,               % Frame around the code
}
\usepackage{float}
\usepackage[]{amsthm} %lets us use \begin{proof}
	\usepackage[]{amssymb} %gives us the character \varnothing
	
	\title{Assignment 3, MATH 4155}
	\author{Zongyi Liu}
	\date{Wed, Feb 25, 2026}
	\begin{document}
		\maketitle
		
		\section{Question 1}
		
		{\color{deepred}\fontspec{Georgia} HW6 Q1 Lyapunov Inequality}
		
		If $X$ is a measurable function on a probability space $(\Omega, \mathcal{F}, \mathbb{P})$ and $0<p<q<\infty$, then
		
		$$
		\|X\|_{p} \leq\|X\|_{q}
		$$
		
		In particular, for $0<p<q<\infty$ we have $\mathbb{L}^{q} \subseteq \mathbb{L}^{p}$ on a probability space. (This kind of inclusion can fail, and rather dramatically, on a general measure space.)
		
		
		
		\textbf{Answer}
		
		\clearpage
		
		\section{Question 2}
		
		{\color{deepred}\fontspec{Georgia} HW6 Q3 An Extended Monotone Convergence Theorem}
		
		Suppose $X_{1} \leq \cdots \leq X_{n} \leq X_{n+1} \leq \cdots$ are integrable random variables, monotonically increasing pointwise to the random variable $X=\lim _{n \rightarrow \infty} X_{n}$.
		
		If $\sup _{n \in \mathbb{N}} \mathbb{E}\left(X_{n}\right)<\infty$, show that $X$ is integrable and that we have
		
		$$
		\lim _{n \rightarrow \infty} \mathbb{E}\left(X_{n}\right)=\mathbb{E}(X)
		$$
		
		\textbf{Answer}
		
			
		
		\clearpage
		
		\section{Question 3}
		
		{\color{deepred}\fontspec{Georgia} HW6 Q4 Riemann-Zeta Calculus} 
		
		For $s>1$ fixed, let $X:\Omega\to\mathbb{N}$ be a random variable with the ``zeta distribution''
		\[
		\mathbb{P}(X=n)=\frac{1}{\zeta(s)\cdot n^{s}}, \qquad n\in\mathbb{N}, 
		\qquad \text{where} \qquad 
		\zeta(s):=\sum_{n\in\mathbb{N}}\frac{1}{n^{s}}
		\]
		is the Riemann zeta-function. For each $m\in\mathbb{N}$, consider the event
		\[
		A_m:=\{\omega\in\Omega: m \text{ is a factor of } X(\omega)\}
		=\{\omega\in\Omega: X(\omega)=m k(\omega), \text{ for some } k(\omega)\in\mathbb{N}\}.
		\]
		
		\begin{enumerate}
			\item[(i)] Show that $\mathbb{P}(A_m)=1/m^{s}$.
			
			\item[(ii)] With $\mathfrak{P}$ denoting the set of prime numbers, conclude from (i) that the events in the collection $\{A_p\}_{p\in\mathfrak{P}}$ are independent,
			
			\item[(iii)] Use (i) and (ii) to derive the celebrated Euler formula
			\[
			\frac{1}{\zeta(s)}=\prod_{p\in\mathfrak{P}}\left(1-\frac{1}{p^{s}}\right).
			\]
		\end{enumerate}
		
		\textbf{Answer}
		
			
		
		\clearpage
		
		\section{Question 4}
		
		{\color{deepred}\fontspec{Georgia} HW6 Q9 Modes of Convergence of Random Variables}
		
		On a suitable probability space, construct
		\begin{enumerate}
			\item[(i)] a sequence of random variables that converges in probability, but not almost surely;
			
			\item[(ii)] a sequence of random variables that converges almost surely, but does not exhibit fast convergence in probability; and
			
			\item[(iii)] a sequence of random variables that converges almost surely, but not in $\mathbb{L}^{1}$.
		\end{enumerate}
		
		(\emph{Hint}: All these sequences can be constructed on the space $\Omega=(0,1]$ with Lebesgue measure $\mathbb{P}$ on its Borel sets)
		
		
		\textbf{Answer} 
		
		
		
		\clearpage
		
		\section{Question 5}
		
		{\color{deepred}\fontspec{Georgia} HW6 Q12 The Concentration of Measure Phenomenon}
			
			\begin{itemize}
				
				\item[(i)] Suppose $X_1,X_2,\cdots$ are square-integrable random variables with
				$m=\mathbb{E}(X_j)$,
				$\mathbb{E}[(X_i-m)(X_j-m)]=0$ for $i\ne j$, and
				$\operatorname{Var}(X_j)=\mathbb{E}[(X_j-m)^2]\le C$
				for all $j$ and some constants $m$ and $C>0$.
				
				Show that, both in $\mathbb{L}^2$ and in probability,
				\[
				\lim_{n\to\infty}\frac{1}{n}\sum_{j=1}^n X_j = m.
				\]
				
				\item[(ii)] Suppose now that $Y_1,Y_2,\cdots$ are independent random variables with
				common distribution uniform on $(-1,1)$.
				Argue that
				\[
				\lim_{n\to\infty}\frac{1}{n}\sum_{j=1}^n Y_j^2=\frac{1}{3}
				\]
				holds in probability. (Does this hold also in a stronger sense?)
				
				\item[(iii)] Consider the cube $\mathbb{K}^n := (-1,1)^n$ in $\mathbb{R}^n$.
				Justify the statement:
				
				\emph{``For $n$ large, most of the volume of this cube is concentrated very
					near the sphere of radius $\sqrt{n/3}$ in $\mathbb{R}^n$.''}
				
				Establish for every $\varepsilon\in(0,1)$ that
				\[
				\lim_{n\to\infty}
				\frac{1}{2^n}\lambda_n\!\left(A_{n,\varepsilon}\cap\mathbb{K}^n\right)=1,
				\]
				where $\lambda_n$ denotes Lebesgue measure in $\mathbb{R}^n$ and
				$\lambda_n(\mathbb{K}^n)=2^n$, and
				\[
				A_{n,\varepsilon}
				=
				\left\{
				y\in\mathbb{R}^n:
				\frac{n}{3}(1-\varepsilon)
				\le
				y_1^2+\cdots+y_n^2
				\le
				\frac{n}{3}(1+\varepsilon)
				\right\}.
				\]
				
			\end{itemize}
		
		\textbf{Answer} 
		
		
	\end{document}
